%!TEX root = ./../../springer-mono-format/book.tex
%!TEX root = ./../../thesis-trier.tex

\IN{$\rho$-automaton}
In general, the equivalence relation~$\eqreach$ induced by the reachability quasiorder is not a $\TS{\mathcal A}$-congruence.
For example, consider the automaton from \cref{prelim:fig:monotonic}
where $1 \eqreach 2$, but $\delta(1, a) \not\eqreach \delta(2, a)$.
An automaton $\mathcal A$ is said to be a \emph{$\rho$-automaton}
if~$\eqreach$ is a $\TS{\mathcal A}$-congruence. 
In this case, the strongly connected components are congruence classes.
Trivially, every strongly connected automaton is a $\rho$-automaton.

\begin{figure}[htb]
	\centering 	
	\begin{subfigure}[b]{0.48\textwidth}
	\centering
		\begin{tikzpicture}[node distance = 5em, on grid, auto, ->]
			\node (q1) [state] {$q_1$};
			\node (q2) [state, above of = q1] {$q_2$};
			\node (q3) [state, right of = q2] {$q_3$};
			\node (q4) [state, left of = q2] {$q_4$};
			
			\path[->] (q1) edge node {$b$} (q2)
			(q2) edge [above] node {$b$} (q4)
			(q3) edge [loop above] node {$a,b$} (q3)
			(q4) edge [loop above] node {$a,b$} (q4)
			(q1) edge [bend left] node {$a$} (q4)
			(q2) edge node {$a$} (q3);
		\end{tikzpicture}
	\caption{Here $q_1 \lereach q_2$, but $\applytrans{\delta}{q_1}{a} \not\lereach \applytrans{\delta}{q_2}{a}$.}
	\label{prelim:fig:rho-aut-left}
	\end{subfigure}\hfill
	\begin{subfigure}[b]{0.5\textwidth}
	\centering
		\begin{tikzpicture}[node distance = 5em, on grid, auto, ->]
			%\tikzstyle{every state}=[inner sep = 5pt, minimum size = 0pt];
			\node (q1) [state] {$q_1$};
			\node (q2) [state, right of = q1] {$q_2$};
			\node (q3) [state, right of = q2] {$q_3$};
			\node (q4) [state, right of = q3] {$q_4$};
			
			\path (q1) edge node {$a$} (q2)
			(q3) edge node {$b$} (q4)
			(q3) edge [loop above] node {$a$} (q3)
			(q4) edge [loop above] node {$a,b$} (q4)
			(q1) edge [bend right, out=320, in = 210, below] node {$b$} (q3)
			(q2) edge [loop above] node {$a,b$} (q2);
		\end{tikzpicture}		
	\caption{Here $q_1 \lereach q_2$, but $\applytrans{\delta}{q_1}{b} \not\lereach \applytrans{\delta}{q_2}{b}$.} 
	\end{subfigure}
	\caption[Example $\rho$-automata]{Two $\rho$-automata where~\cref{constr_sync:eqn:reach_compatible}
		is not true and $\delta(q_1, a) \not\lereach q_1$.}
	\label{contr_input:fig:rho_aut}
\end{figure}

First, we give two sufficient conditions to identify $\rho$-automata.

\begin{proposition}\label{prelim:prop:rho-automata}
	Let $\semiautomat{A}{Q}{\Sigma}{\delta}$.
	
	\begin{enumerate}
	\item \label{prelim:prop:rho-automata:1} Suppose,  for all $p,q \in Q$ and $a \in \Sigma$, we have
	\begin{equation}\label{constr_sync:eqn:reach_compatible}
		p \lereach q \Rightarrow \applytrans{\delta}{p}{a} \lereach \applytrans{\delta}{q}{a}.
	\end{equation}
	Then $\mathcal A$ is a $\rho$-automaton.
	
	\item \label{prelim:prop:rho-automata:2} Suppose, for all $q\in Q$ and $a \in \Sigma$,
	 we have $\applytrans{\delta}{q}{a} \lereach q$.
	 Then $\mathcal A$ 
	 fulfills \cref{constr_sync:eqn:reach_compatible} and so it is a $\rho$-automaton.
	\end{enumerate}
\end{proposition}
\begin{proof}
\ref{prelim:prop:rho-automata:1}: Let $a \in \Sigma$ and $p,q \in Q$. If $p \eqreach q$, 
then $\applytrans{\delta}{p}{a} \eqreach \applytrans{\delta}{q}{a}$.	
Hence Equation~\eqref{prelim:eqn:TS-congruence} 
holds and so $\eqreach$ is a $\TS{\mathcal A}$-congruence.
	
\ref{prelim:prop:rho-automata:2}: Let $p,q \in Q$ with $p \lereach q$ and $a \in \Sigma$.
Then, by assumption and~\cref{prelim:prop:reachability-finest},
we have $\applytrans{\delta}{p}{a} \lereach p \lereach q \lereach \applytrans{\delta}{q}{a}$.
Whence~\cref{constr_sync:eqn:reach_compatible}
is true and the claim follows from~\cref{prelim:prop:rho-automata:1}.
\end{proof}

Figure~\ref{contr_input:fig:rho_aut} shows
two $\rho$-automata where both conditions
from~\cref{prelim:prop:rho-automata} do not hold,
so that they give proper subclasses of $\rho$-automata.
Also, the classes of partially ordered
and permutation automata form proper subclasses
of $\rho$-automata.

\begin{proposition}
	\label{constr_sync:prop:pto_comm_rho_automata}
	Every partially ordered automaton is a $\rho$-automaton
	and every permutation automaton is a $\rho$-automaton.
\end{proposition}
\begin{proof}
	For a partially ordered automaton $\mathcal A$, the relation
	$\eqreach$ equals the equality relation, which is a trivial $\TS{\mathcal A}$-congruence. 
    By \cref{prelim:subclasses:perm:reach}
	and~\cref{prelim:prop:rho-automata:2} of \cref{prelim:prop:rho-automata}
	we find that permutation automata are $\rho$-automata.
\end{proof}	

The reachability order on $\mathcal A / {\eqreach}$
equals the order induced on the $\eqreach$-classes by the reachability order on $\mathcal A$.
Since this is a partial order, we find that $\mathcal A / {\eqreach}$
is a partially ordered automaton.
See \cref{fig:ex_comm_aut}
for an example. 

Moreover, constructing $\mathcal A / {\eqreach}$ amounts to contracting the
strongly connected components of $\mathcal A$.\footnote{In graph theory, the
	graph obtained by contracting each strongly connected component to a single
	vertex is called the \emph{condensation} of the graph.} With Tarjan's
algorithm~\cite{DBLP:journals/siamcomp/Tarjan72} this can be done in time  $O(|Q| + |\Sigma||Q|)$.
We thus obtain the following statement.

\begin{figure}[htb]
     \centering
    \scalebox{1.0}{
% diamand in der mitte
% \begin{tikzpicture}
%  \tikzset{every state/.style={minimum size=1pt},>=stealth'}
%  \tikzstyle{dotted}=                  [dash pattern=on \pgflinewidth off 2pt]
   
 \begin{tikzpicture}[
 	>=latex',
 	shorten >=1pt,
 	node distance=6em,
 	on grid,
 	auto
 	]
 	\tikzset{every state/.style={minimum size=1pt},
 		mybox/.style={
 			draw=red,
 			dashed,
 			rounded corners,
 			inner sep=13pt
 		}
 	}
 	
 	%---------------- Left automaton ----------------%
 	
 	\node[state] (2) {};
 	\node[state] (3) [below left=of 2] {};
 	\node[state] (4) [below right=of 2] {};
 	\node[state] (5) [left=of 3] {};
 	\node[state] (6) [below right=of 3] {};
 	\node[state] (8) [above right=of 4] {};
 	\node[state] (9) [above=of 8] {};
 	
 	\path[->]
 	(2) edge node {$a$} (3)
 	(3) edge[bend left] node {$a$} (5)
 	(5) edge[bend left] node {$a$} (3)
 	(3) edge node {$b$} (6)
 	(6) edge[loop below] node {$a,b$} (6)
 	(5) edge[bend right=70] node {$b$} (6);
 	
 	\path[->]
 	(2) edge node {$b$} (4)
 	(4) edge[loop right] node {$b$} (4)
 	(4) edge node {$a$} (6);
 	
 	\path[->]
 	(8) edge[bend left] node {$b$} (9)
 	(9) edge[bend left] node {$b$} (8)
 	(8) edge node {$a$} (4)
 	(9) edge[bend right] node {$a$} (4);
 	
 	%---------------- Arrow ----------------%
 	
 	\node (P) [right=6em of 4] {\LARGE $\Rightarrow$};
 	
 	%---------------- Right automaton ----------------%
 	
 	\node[state] (S1) [right=3em of P] {};
 	\node[state] (S2) [above right=of S1] {};
 	\node[state] (S3) [below right=of S2] {};
 	\node[state] (S4) [below left=of S3] {};
 	\node[state] (S5) [above right=of S3] {};
 	
 	\path[->]
 	(S2) edge node {$a$} (S1)
 	(S2) edge node {$b$} (S3)
 	(S3) edge node {$a$} (S4)
 	(S1) edge node {$b$} (S4)
 	(S5) edge node {$a$} (S3);
 	
 	\path[->]
 	(S1) edge[loop above] node {$a$} (S1)
 	(S3) edge[loop right] node {$b$} (S3)
 	(S5) edge[loop above] node {$a,b$} (S5)
 	(S4) edge[loop below] node {$a,b$} (S4);
 	
 	%---------------- Background rectangles ----------------%
 	
 	\begin{pgfonlayer}{background}
 		
 		% States 3 and 5
 		\node[mybox, fit=(3)(5)] {};
 		
 		% State 6
 		\node[
 		mybox,
 		fit=(6), inner sep=8
 		] {};
 		
 		% State 2
 		\node[
 		mybox,
 		fit=(2), inner sep=8
 		] {};
 		
 		% State 4
 		\node[
 		mybox,
 		fit=(4), inner sep=8
 		] {};
 		
 		% States 8 and 9
 		\node[
 		mybox,
 		fit=(8)(9)
 		] {};
 		
 	\end{pgfonlayer}
 	
 \end{tikzpicture}
}
  \caption[Example quotient automaton of \texorpdfstring{$\rho$}{rho}-automaton]{
  	Left: A $\rho$-automaton $\mathcal A$ (that is commutative and synchronizing). The strongly connected components of $\mathcal A$
  	are surrounded by boxes. Right: The automaton $\mathcal A / {\eqreach}$.}
  %    have been framed by red boxes.}
  \label{fig:ex_comm_aut}
\end{figure}





%With Theorem~\ref{sync_max:thm:com_waa_same_sync_words}, 
%a commutative automaton with $n$
%states can only be sync-maximal when $2^n - n \le n$.
%This leaves only the possibilities $n \in \{1,2\}$.
%As sync-maximality implies that it is synchronizing \todo{das irgendwo hinschreiben als claim oder so?}
%the claim follows with Proposition~\ref{prop:two_state_comm_sync}.
%Conversely, that the two types of automata are commutative
%and sync-maximal is easily seen.



%\begin{toappendix2} 
%
%\begin{figure}[htb]
%     \centering
%    \scalebox{.7}{
%% diamand in der mitte
%% \begin{tikzpicture}
%%  \tikzset{every state/.style={minimum size=1pt},>=stealth'}
%%  \tikzstyle{dotted}=                  [dash pattern=on \pgflinewidth off 2pt]
%   
%  \begin{tikzpicture}[>=latex',shorten >=1pt,node distance=2cm,on grid,auto]
% \tikzset{every state/.style={minimum size=1pt}}
%    \node[state] (10) at (11,2) {};
%   \node[state] (11) at (14,1) {}; %[below right of=10] {};
%   \node[state] (12) at (10.5,-0.5) {}; % [below of=11] {};
%   \node[state] (13) [below of=12] {};
%   \node[state] (14) [right of=13] {};
%   \node[state] (15) [left of=11] {};
%   \node[state] (16) [right of=12] {};
%   
%   \path[->] (10) edge [bend left] node {$a$} (11)
%             (11) edge node {$b$} (16)
%             (10) edge [left] node {$b$} (12)
%            % (11) edge [loop left] node {$a$} (11)
%             (12) edge [left] node {$b$} (13)
%             (13) edge [loop below] node {$b$} (12)
%             (14) edge [loop below] node {$b$} (14)
%             (13) edge [bend left] node {$a$} (14)
%             (14) edge [bend left] node {$a$} (13);
%    
% % \path[->] (12) edge [loop right] node {$a$} (12);
% 
%   \path[->] (11) edge [bend left] node {$a$} (15)
%             (15) edge [bend left] node {$a$} (11);
%             
%   \path[->] (16) edge [bend left] node {$a$} (12)
%             (12) edge [bend left] node {$a$} (16);
%             
%   \path[->] (15) edge [left,pos=.3] node {$b$} (12);
%   \path[->] (16) edge node {$b$} (14);
%   
%\end{tikzpicture}
%}
%   
%   \caption{A commutative automaton that is not synchronizing.}
%  \label{fig:ex_comm_aut2}
%\end{figure}
%\end{toappendix2}

\begin{proposition}
	\label{constr_input:prop:factor_rho_aut_pto}
	Let $\semiautomat{A}{Q}{\Sigma}{\delta}$ be a $\rho$-automaton.
	Then $\mathcal A / {\eqreach}$ is a partially ordered automaton
	computable in $O(|Q| + |\Sigma||Q|)$.
\end{proposition}

Next, we show that for $\rho$-automata with a trap state, the set of synchronizing words is preserved when considering the quotient automaton $\mathcal A / {\eqreach}$.

\begin{proposition}
	\label{constr_input:prop:rho_trap_sync_factor}
	Let $\semiautomat{A}{Q}{\Sigma}{\delta}$ be a $\rho$-automaton
	with a trap state.
	Then $\Syn(\mathcal A) = \Syn(\mathcal A / \eqreach)$.
\end{proposition}
\begin{proof}
	For the reachability order,
	we have $[t]_{\eqreach} = \{t\}$ for each trap state~$t$.
	Thus the claim follows from \cref{prelim:syncaut:cor:Syncongr}.	
\end{proof}

The assumption that $\mathcal A$
has a trap state is crucial. For example,
for the \Cerny family~$\mathcal C_n$ 
(Figure~\ref{intro:fig:cerny})
we have $\Syn(\mathcal C_n) \ne \Sigma^*$
but $\Syn(\mathcal C_n / {\eqreach}) = \Sigma^*$.
%(the latter equality, in fact, holds for every strongly
%connected automaton).

\begin{corollary} \label{constr_input:prop:rho_trap_resetlength}
	Every synchronizing $\rho$-automaton $\mathcal A$ with $n$ states and a trap state
	has a synchronizing word of length at most $n-1$,
	and this bound is tight.
\end{corollary}
\begin{proof}
	Since $\mathcal A / {\eqreach}$ is partially ordered, its strongly connected
	components form blocks. Hence, by~\cref{prelim:thm:SCCblocks-rank_r}
	with $r = 1$, there exists a synchronizing word of length at most $n-1$
	in $\Syn(\mathcal A / {\eqreach})$.
	The result now follows from
	$\Syn(\mathcal A / {\eqreach}) = \Syn(\mathcal A)$,
	which holds by~\cref{constr_input:prop:rho_trap_sync_factor}.
\end{proof}

\begin{corollary} \label{constr_input:prop:rho_extremal_sync_then_pto}
	A synchronizing $n$-state $\rho$-automaton $\mathcal A$ with a trap state all of whose
	synchronizing words have length at least $n - 1$
	is partially ordered.
\end{corollary} 
\begin{proof}
	If it is not partially ordered, then $\mathcal A / {\eqreach}$
	has strictly less states. Since it has a trap state,
	it has a synchronizing word
	of length less than $n - 1$ by \cref{constr_input:prop:rho_trap_resetlength},
	which also synchronizes $\mathcal A$
	by \cref{constr_input:prop:rho_trap_sync_factor}.
	Thus, $\mathcal A$ must be partially ordered.	
\end{proof}